\documentclass[9pt,t]{beamer}
% =====================================================================
% finance-beamer.tex — shared Beamer style for the LLM-in-Finance course
% =====================================================================
% Reusable slide style. Each deck \input's this immediately after
% \documentclass{beamer}. It provides:
%   * the Madrid theme + book colour palette (matches the printed book)
%   * two book-style framing blocks:
%       \begin{bigpicture} ... \end{bigpicture}   ("the bigger picture")
%       \begin{underhood}   ... \end{underhood}    ("under the hood")
%     These mirror the book's `context` and `deepdive` tcolorboxes so the
%     same intuition-vs-mechanism scaffold the reader meets in the book
%     reappears on the slides.
%   * a Python listings style for code frames
%   * appendix conventions: \appendixstart drops a divider and switches
%     to non-numbered backup frames so "extra / less central" material
%     lives after the main talk without inflating the slide count.
%
% Usage:
%   \documentclass{beamer}
%   % =====================================================================
% finance-beamer.tex — shared Beamer style for the LLM-in-Finance course
% =====================================================================
% Reusable slide style. Each deck \input's this immediately after
% \documentclass{beamer}. It provides:
%   * the Madrid theme + book colour palette (matches the printed book)
%   * two book-style framing blocks:
%       \begin{bigpicture} ... \end{bigpicture}   ("the bigger picture")
%       \begin{underhood}   ... \end{underhood}    ("under the hood")
%     These mirror the book's `context` and `deepdive` tcolorboxes so the
%     same intuition-vs-mechanism scaffold the reader meets in the book
%     reappears on the slides.
%   * a Python listings style for code frames
%   * appendix conventions: \appendixstart drops a divider and switches
%     to non-numbered backup frames so "extra / less central" material
%     lives after the main talk without inflating the slide count.
%
% Usage:
%   \documentclass{beamer}
%   % =====================================================================
% finance-beamer.tex — shared Beamer style for the LLM-in-Finance course
% =====================================================================
% Reusable slide style. Each deck \input's this immediately after
% \documentclass{beamer}. It provides:
%   * the Madrid theme + book colour palette (matches the printed book)
%   * two book-style framing blocks:
%       \begin{bigpicture} ... \end{bigpicture}   ("the bigger picture")
%       \begin{underhood}   ... \end{underhood}    ("under the hood")
%     These mirror the book's `context` and `deepdive` tcolorboxes so the
%     same intuition-vs-mechanism scaffold the reader meets in the book
%     reappears on the slides.
%   * a Python listings style for code frames
%   * appendix conventions: \appendixstart drops a divider and switches
%     to non-numbered backup frames so "extra / less central" material
%     lives after the main talk without inflating the slide count.
%
% Usage:
%   \documentclass{beamer}
%   \input{../_style/finance-beamer.tex}
%   \title{...}\subtitle{...}\author{...}\institute{...}\date{\today}
%   \begin{document} ... \end{document}
% =====================================================================

\usetheme{Madrid}
\usecolortheme{whale}
\setbeamertemplate{navigation symbols}{}
\setbeamertemplate{caption}[numbered]
\setbeamercovered{transparent}

% --- Density controls: keep dense, equation-heavy frames on one slide ---
% Slightly smaller body text + tighter lists/blocks. Frame titles and the
% Madrid chrome keep their normal size, so the deck still reads as Madrid,
% just with more vertical room for content.
\setbeamerfont{normal text}{size=\small}
\setbeamertemplate{itemize/enumerate body begin}{\small}
\setbeamertemplate{itemize/enumerate subbody begin}{\footnotesize}
% Tighter block spacing.
\setbeamertemplate{block begin}{%
  \par\vskip2pt%
  \begin{beamercolorbox}[colsep*=.75ex,leftskip=1ex,rightskip=1ex]{block title}%
    \usebeamerfont*{block title}\insertblocktitle%
  \end{beamercolorbox}%
  {\parskip0pt\vskip-1pt}%
  \begin{beamercolorbox}[colsep*=.75ex,vmode,leftskip=1ex,rightskip=1ex]{block body}%
  \vskip1pt}
\setbeamertemplate{block end}{\end{beamercolorbox}\vskip2pt}

\usepackage{amsmath, amssymb}
\usepackage{booktabs}
\usepackage{xcolor}
\usepackage{listings}
\usepackage[skins, breakable]{tcolorbox}

% --- Book colour palette (kept in sync with book/preamble.tex) ---
\definecolor{linkblue}{RGB}{14, 75, 140}
\definecolor{deepdivebg}{RGB}{235, 244, 255}
\definecolor{narrativebg}{RGB}{255, 252, 240}
\definecolor{narrativerule}{RGB}{180, 130, 40}
\definecolor{steelblue}{RGB}{70, 130, 180}
\definecolor{forestgreen}{RGB}{34, 139, 34}
\definecolor{crimson}{RGB}{220, 20, 60}
\definecolor{darkorange}{RGB}{255, 140, 0}
\definecolor{codebg}{RGB}{248, 248, 248}
\definecolor{coderule}{RGB}{210, 210, 210}
\definecolor{keywordblue}{RGB}{0, 100, 180}
\definecolor{commentgray}{RGB}{120, 120, 120}
\definecolor{stringgreen}{RGB}{30, 130, 30}

% Tie the Madrid structural colour to the book's link blue.
\setbeamercolor{structure}{fg=linkblue}
\setbeamercolor{title}{fg=white}
\setbeamercolor{frametitle}{fg=white}

% --- "the bigger picture" : narrative / intuition framing -----------
\newtcolorbox{bigpicture}{%
  enhanced, breakable, boxsep=2pt,
  colback  = narrativebg,
  colframe = narrativerule,
  boxrule  = 0pt, toprule = 1.4pt, bottomrule = 0.4pt,
  leftrule = 0pt, rightrule = 0pt, arc = 0pt,
  left = 6pt, right = 6pt, top = 4pt, bottom = 5pt,
  fonttitle = \footnotesize\scshape\color{narrativerule},
  title = {the bigger picture}, attach title to upper = {\par\smallskip},
}

% --- "under the hood" : the mechanism / technical detail ------------
\newtcolorbox{underhood}{%
  enhanced, breakable, boxsep=2pt,
  colback  = deepdivebg,
  colframe = linkblue,
  boxrule  = 0pt, toprule = 1.4pt, bottomrule = 0.4pt,
  leftrule = 0pt, rightrule = 0pt, arc = 0pt,
  left = 6pt, right = 6pt, top = 4pt, bottom = 5pt,
  fonttitle = \footnotesize\scshape\color{linkblue},
  title = {under the hood}, attach title to upper = {\par\smallskip},
}

% --- Python code style ---------------------------------------------
\lstset{
  basicstyle    = \ttfamily\footnotesize,
  keywordstyle  = \color{keywordblue}\bfseries,
  commentstyle  = \color{commentgray}\itshape,
  stringstyle   = \color{stringgreen},
  showstringspaces = false,
  language      = Python,
  breaklines    = true,
  backgroundcolor = \color{codebg},
  frame         = single,
  rulecolor     = \color{coderule},
  numbers       = none,
  aboveskip     = 4pt, belowskip = 4pt,
}

% --- Appendix conventions ------------------------------------------
% \appendixstart{Title} : begins the "extra material" appendix. Frames
% after it are not counted toward the main deck's frame total, keeping
% the core talk lean while preserving the deeper / optional content.
\newcommand{\appendixstart}[1]{%
  \appendix
  \setbeamertemplate{frametitle continuation}{}%
  \begin{frame}[noframenumbering]
    \begin{center}
      {\Large\scshape\color{linkblue} Appendix}\\[4pt]
      {\large #1}\\[10pt]
      {\footnotesize Extra / optional material — beyond the core lecture.}
    \end{center}
  \end{frame}
}
% Use \begin{frame}[noframenumbering]{...} for each appendix frame.

%   \title{...}\subtitle{...}\author{...}\institute{...}\date{\today}
%   \begin{document} ... \end{document}
% =====================================================================

\usetheme{Madrid}
\usecolortheme{whale}
\setbeamertemplate{navigation symbols}{}
\setbeamertemplate{caption}[numbered]
\setbeamercovered{transparent}

% --- Density controls: keep dense, equation-heavy frames on one slide ---
% Slightly smaller body text + tighter lists/blocks. Frame titles and the
% Madrid chrome keep their normal size, so the deck still reads as Madrid,
% just with more vertical room for content.
\setbeamerfont{normal text}{size=\small}
\setbeamertemplate{itemize/enumerate body begin}{\small}
\setbeamertemplate{itemize/enumerate subbody begin}{\footnotesize}
% Tighter block spacing.
\setbeamertemplate{block begin}{%
  \par\vskip2pt%
  \begin{beamercolorbox}[colsep*=.75ex,leftskip=1ex,rightskip=1ex]{block title}%
    \usebeamerfont*{block title}\insertblocktitle%
  \end{beamercolorbox}%
  {\parskip0pt\vskip-1pt}%
  \begin{beamercolorbox}[colsep*=.75ex,vmode,leftskip=1ex,rightskip=1ex]{block body}%
  \vskip1pt}
\setbeamertemplate{block end}{\end{beamercolorbox}\vskip2pt}

\usepackage{amsmath, amssymb}
\usepackage{booktabs}
\usepackage{xcolor}
\usepackage{listings}
\usepackage[skins, breakable]{tcolorbox}

% --- Book colour palette (kept in sync with book/preamble.tex) ---
\definecolor{linkblue}{RGB}{14, 75, 140}
\definecolor{deepdivebg}{RGB}{235, 244, 255}
\definecolor{narrativebg}{RGB}{255, 252, 240}
\definecolor{narrativerule}{RGB}{180, 130, 40}
\definecolor{steelblue}{RGB}{70, 130, 180}
\definecolor{forestgreen}{RGB}{34, 139, 34}
\definecolor{crimson}{RGB}{220, 20, 60}
\definecolor{darkorange}{RGB}{255, 140, 0}
\definecolor{codebg}{RGB}{248, 248, 248}
\definecolor{coderule}{RGB}{210, 210, 210}
\definecolor{keywordblue}{RGB}{0, 100, 180}
\definecolor{commentgray}{RGB}{120, 120, 120}
\definecolor{stringgreen}{RGB}{30, 130, 30}

% Tie the Madrid structural colour to the book's link blue.
\setbeamercolor{structure}{fg=linkblue}
\setbeamercolor{title}{fg=white}
\setbeamercolor{frametitle}{fg=white}

% --- "the bigger picture" : narrative / intuition framing -----------
\newtcolorbox{bigpicture}{%
  enhanced, breakable, boxsep=2pt,
  colback  = narrativebg,
  colframe = narrativerule,
  boxrule  = 0pt, toprule = 1.4pt, bottomrule = 0.4pt,
  leftrule = 0pt, rightrule = 0pt, arc = 0pt,
  left = 6pt, right = 6pt, top = 4pt, bottom = 5pt,
  fonttitle = \footnotesize\scshape\color{narrativerule},
  title = {the bigger picture}, attach title to upper = {\par\smallskip},
}

% --- "under the hood" : the mechanism / technical detail ------------
\newtcolorbox{underhood}{%
  enhanced, breakable, boxsep=2pt,
  colback  = deepdivebg,
  colframe = linkblue,
  boxrule  = 0pt, toprule = 1.4pt, bottomrule = 0.4pt,
  leftrule = 0pt, rightrule = 0pt, arc = 0pt,
  left = 6pt, right = 6pt, top = 4pt, bottom = 5pt,
  fonttitle = \footnotesize\scshape\color{linkblue},
  title = {under the hood}, attach title to upper = {\par\smallskip},
}

% --- Python code style ---------------------------------------------
\lstset{
  basicstyle    = \ttfamily\footnotesize,
  keywordstyle  = \color{keywordblue}\bfseries,
  commentstyle  = \color{commentgray}\itshape,
  stringstyle   = \color{stringgreen},
  showstringspaces = false,
  language      = Python,
  breaklines    = true,
  backgroundcolor = \color{codebg},
  frame         = single,
  rulecolor     = \color{coderule},
  numbers       = none,
  aboveskip     = 4pt, belowskip = 4pt,
}

% --- Appendix conventions ------------------------------------------
% \appendixstart{Title} : begins the "extra material" appendix. Frames
% after it are not counted toward the main deck's frame total, keeping
% the core talk lean while preserving the deeper / optional content.
\newcommand{\appendixstart}[1]{%
  \appendix
  \setbeamertemplate{frametitle continuation}{}%
  \begin{frame}[noframenumbering]
    \begin{center}
      {\Large\scshape\color{linkblue} Appendix}\\[4pt]
      {\large #1}\\[10pt]
      {\footnotesize Extra / optional material — beyond the core lecture.}
    \end{center}
  \end{frame}
}
% Use \begin{frame}[noframenumbering]{...} for each appendix frame.

%   \title{...}\subtitle{...}\author{...}\institute{...}\date{\today}
%   \begin{document} ... \end{document}
% =====================================================================

\usetheme{Madrid}
\usecolortheme{whale}
\setbeamertemplate{navigation symbols}{}
\setbeamertemplate{caption}[numbered]
\setbeamercovered{transparent}

% --- Density controls: keep dense, equation-heavy frames on one slide ---
% Slightly smaller body text + tighter lists/blocks. Frame titles and the
% Madrid chrome keep their normal size, so the deck still reads as Madrid,
% just with more vertical room for content.
\setbeamerfont{normal text}{size=\small}
\setbeamertemplate{itemize/enumerate body begin}{\small}
\setbeamertemplate{itemize/enumerate subbody begin}{\footnotesize}
% Tighter block spacing.
\setbeamertemplate{block begin}{%
  \par\vskip2pt%
  \begin{beamercolorbox}[colsep*=.75ex,leftskip=1ex,rightskip=1ex]{block title}%
    \usebeamerfont*{block title}\insertblocktitle%
  \end{beamercolorbox}%
  {\parskip0pt\vskip-1pt}%
  \begin{beamercolorbox}[colsep*=.75ex,vmode,leftskip=1ex,rightskip=1ex]{block body}%
  \vskip1pt}
\setbeamertemplate{block end}{\end{beamercolorbox}\vskip2pt}

\usepackage{amsmath, amssymb}
\usepackage{booktabs}
\usepackage{xcolor}
\usepackage{listings}
\usepackage[skins, breakable]{tcolorbox}

% --- Book colour palette (kept in sync with book/preamble.tex) ---
\definecolor{linkblue}{RGB}{14, 75, 140}
\definecolor{deepdivebg}{RGB}{235, 244, 255}
\definecolor{narrativebg}{RGB}{255, 252, 240}
\definecolor{narrativerule}{RGB}{180, 130, 40}
\definecolor{steelblue}{RGB}{70, 130, 180}
\definecolor{forestgreen}{RGB}{34, 139, 34}
\definecolor{crimson}{RGB}{220, 20, 60}
\definecolor{darkorange}{RGB}{255, 140, 0}
\definecolor{codebg}{RGB}{248, 248, 248}
\definecolor{coderule}{RGB}{210, 210, 210}
\definecolor{keywordblue}{RGB}{0, 100, 180}
\definecolor{commentgray}{RGB}{120, 120, 120}
\definecolor{stringgreen}{RGB}{30, 130, 30}

% Tie the Madrid structural colour to the book's link blue.
\setbeamercolor{structure}{fg=linkblue}
\setbeamercolor{title}{fg=white}
\setbeamercolor{frametitle}{fg=white}

% --- "the bigger picture" : narrative / intuition framing -----------
\newtcolorbox{bigpicture}{%
  enhanced, breakable, boxsep=2pt,
  colback  = narrativebg,
  colframe = narrativerule,
  boxrule  = 0pt, toprule = 1.4pt, bottomrule = 0.4pt,
  leftrule = 0pt, rightrule = 0pt, arc = 0pt,
  left = 6pt, right = 6pt, top = 4pt, bottom = 5pt,
  fonttitle = \footnotesize\scshape\color{narrativerule},
  title = {the bigger picture}, attach title to upper = {\par\smallskip},
}

% --- "under the hood" : the mechanism / technical detail ------------
\newtcolorbox{underhood}{%
  enhanced, breakable, boxsep=2pt,
  colback  = deepdivebg,
  colframe = linkblue,
  boxrule  = 0pt, toprule = 1.4pt, bottomrule = 0.4pt,
  leftrule = 0pt, rightrule = 0pt, arc = 0pt,
  left = 6pt, right = 6pt, top = 4pt, bottom = 5pt,
  fonttitle = \footnotesize\scshape\color{linkblue},
  title = {under the hood}, attach title to upper = {\par\smallskip},
}

% --- Python code style ---------------------------------------------
\lstset{
  basicstyle    = \ttfamily\footnotesize,
  keywordstyle  = \color{keywordblue}\bfseries,
  commentstyle  = \color{commentgray}\itshape,
  stringstyle   = \color{stringgreen},
  showstringspaces = false,
  language      = Python,
  breaklines    = true,
  backgroundcolor = \color{codebg},
  frame         = single,
  rulecolor     = \color{coderule},
  numbers       = none,
  aboveskip     = 4pt, belowskip = 4pt,
}

% --- Appendix conventions ------------------------------------------
% \appendixstart{Title} : begins the "extra material" appendix. Frames
% after it are not counted toward the main deck's frame total, keeping
% the core talk lean while preserving the deeper / optional content.
\newcommand{\appendixstart}[1]{%
  \appendix
  \setbeamertemplate{frametitle continuation}{}%
  \begin{frame}[noframenumbering]
    \begin{center}
      {\Large\scshape\color{linkblue} Appendix}\\[4pt]
      {\large #1}\\[10pt]
      {\footnotesize Extra / optional material — beyond the core lecture.}
    \end{center}
  \end{frame}
}
% Use \begin{frame}[noframenumbering]{...} for each appendix frame.


\title{Loops, Goals, and Iterations}
\subtitle{Lecture 17 --- Agents, Skills, and Hooks}
\author{Juan F.\ Imbet}
\institute{Paris Dauphine -- PSL University}
\date{\today}

\begin{document}

\begin{frame}\titlepage\end{frame}

\begin{frame}{Outline}\tableofcontents\end{frame}

% =====================================================================
\section{From Linear Prompts to Iterative Systems}
% =====================================================================

\begin{frame}{The Prompt Is Not the Unit of Design --- the Loop Is}

\begin{bigpicture}
Earlier lectures treated the LLM as a \emph{function}: prompt in, completion out.
This lecture's thesis: almost every consequential use of an LLM in finance is a
\emph{process} --- a sequence of model calls, tool calls, validations, and
revisions around a goal the first call never achieves.
\end{bigpicture}

\medskip
\textbf{The methodology capstone of the course.} We study the \alert{design of
iterative agentic systems}, not any single tool. Concrete tooling (Claude Code,
Codex, the SDKs, Git) lives in Appendices A--F.

\medskip
\begin{block}{What you will leave with}
A vocabulary --- \emph{goals, agents, skills, hooks, stopping rules} --- and two
worked loops: a transparent root-finder and a nine-agent credit-stress loop.
\end{block}
\end{frame}

\begin{frame}{Why One-Shot Prompting Is Fragile}

A one-shot prompt is an \alert{open-loop controller}: output consumed downstream
with \emph{no error-correction channel}. The model never learns whether it was
right. Fragility compounds along three axes:

\medskip
\begin{enumerate}
  \item \textbf{Errors propagate without attenuation.} One forward pass to parse,
        extract, compute, narrate $\Rightarrow$ an early error corrupts all that
        follows, unnoticed.
  \item \textbf{The spec exceeds one pass.} Interacting constraints (turnover,
        sector caps, tracking error, tax lots) --- joint satisfaction falls
        roughly \emph{multiplicatively}.
  \item \textbf{The model cannot see the environment.} It guesses the schema, the
        rate, whether a file exists --- acting on a \emph{hallucinated} world.
\end{enumerate}

\medskip
\begin{bigpicture}
No analyst writes a DCF top-to-bottom and submits it. They sketch, sanity-check
the terminal value, walk back an over-ceiling margin, re-run, circulate. The
expert workflow is intrinsically \emph{iterative}.
\end{bigpicture}
\end{frame}

\begin{frame}{The Goal-Directed Loop}

\begin{block}{Five phases --- repeat until $\mathrm{ok}(s)$ holds or a stopping rule fires}
\textbf{Observe} $\to$ \textbf{Propose} $\to$ \textbf{Act} $\to$
\textbf{Evaluate} $\to$ \textbf{Revise}
\end{block}

\smallskip
\begin{itemize}
  \item \textbf{Observe} state $s_t$ (files, tests, data, prior outputs).
  \item \textbf{Propose} action $a_t$ toward goal $G$, conditioned on $s_t$.
  \item \textbf{Act}: execute $a_t \Rightarrow s_{t+1}$.
  \item \textbf{Evaluate}: feedback $e_{t+1}$ from checks; does $\mathrm{ok}(s_{t+1})$ hold?
  \item \textbf{Revise}: fold $e_{t+1}$ into context, return to Observe.
\end{itemize}

\medskip
The ordering is a dependency chain. \alert{Revise feeds the concrete evidence of
what went wrong back into the next Propose} --- the error-correction channel
one-shot prompting lacks.

\medskip
\begin{alertblock}{The Evaluate phase carries most of the value}
A loop is only as good as its feedback. A weak $\mathrm{ok}(\cdot)$ converges
happily to something useless. Design the feedback like you design the prompt.
\end{alertblock}
\end{frame}

\begin{frame}{Loops vs.\ Repetition --- and When \emph{Not} to Loop}

\begin{underhood}
\textbf{Genuine iteration:} $a_{t+1} = \pi(s_{t+1}, e_{t+1}, h_t)$ --- the next
action depends on the \emph{evaluated feedback}. If $a_{t+1}\perp e_{t+1}$, it is
\alert{mere repetition} (sampling), not iteration.
\end{underhood}

\medskip
\textbf{Anti-pattern --- spurious repetition:} re-prompt ``fix the failing test''
without reading whether it now passes. Same number of calls; \emph{no information
carried forward}; partial edits may leave the code worse.

\medskip
\textbf{Decision heuristic --- loop only if:} output is mechanically verifiable;
reasoning varies with input; first attempt likely wrong \emph{and} costly; no
irreversible harm before a human gate.

\medskip
\begin{tabular}{@{}lll@{}}
\toprule
\textbf{Task} & \textbf{Right tool} & \textbf{Why} \\
\midrule
Bond yield to maturity & Pure function & Deterministic \\
Daily price-file download & Cron job & Scheduled, no reasoning \\
Figures from a novel filing & Loop + validator & Checkable output \\
Harden a valuation model & Agentic loop & Costly silent errors \\
\bottomrule
\end{tabular}
\end{frame}

% =====================================================================
\section{Goals as Control Objects}
% =====================================================================

\begin{frame}{Goals as Control Objects: Concrete, Checkable, Decomposable}

The goal is the loop's \alert{setpoint} --- it governs behavior, decides when to
stop, and decides what is acceptable. Three properties, one per phase:

\medskip
\begin{itemize}
  \item \textbf{Concrete} (for \emph{Act}) --- names the artifact.
        ``Improve the model'' $\to$ ``raise OOS AUC of \texttt{credit.py},
        measured by \texttt{backtest.py}, to $\geq 0.78$.''
  \item \textbf{Checkable} (for \emph{Evaluate}) --- supplies the predicate
        $\mathrm{ok}(\cdot)$: ``every rubric dimension $\geq 90$ via
        \texttt{/score-content}.''
  \item \textbf{Decomposable} (for \emph{Revise}) --- splits into independently
        checkable subgoals; localizes failures and bounds context.
\end{itemize}

\medskip
\begin{block}{A complete goal specifies \emph{three} distinct things}
\textbf{Success criterion} (stop because we won) $\cdot$
\textbf{Stopping rule} (stop anyway: budget, no-progress) $\cdot$
\textbf{Anticipated failure modes} (named in advance, each becomes a check).
The common defect: specifying only the first.
\end{block}
\end{frame}

\begin{frame}[fragile]{Write the Goal as a File --- Ideally Executable}

A goal in a chat prompt vanishes when context rolls over. A goal in a \emph{file}
is durable, auditable, version-controlled, and --- when executable --- usable
\alert{directly as the Evaluate phase}.

\begin{lstlisting}[language={}]
# TASK: Harden the FCF valuation script
## Artifact
valuation/dcf.py  (value_company(ticker) -> float)
## Success criteria
- pytest tests/test_dcf.py exits 0
- Coverage of dcf.py >= 90%
- No division-by-zero on the 12 edge-case tickers
## Stopping rule
- Halt after 6 iterations, or 2 with no new passing tests.
## Anticipated failure modes
- Passing tests by catching all exceptions and returning 0
  -> forbidden; every value finite & positive or raise ValueError.
\end{lstlisting}

\smallskip
The decisive move: make ``done'' and ``acceptable'' \emph{the same code} --- a
\texttt{validate.sh} that exits 0 iff the goal is met, wired into Evaluate and
enforced before commit.
\end{frame}

\begin{frame}{Bad Goals vs.\ Good Goals --- a Diagnostic}

\begin{tabular}{@{}p{0.30\linewidth}p{0.40\linewidth}p{0.22\linewidth}@{}}
\toprule
\textbf{Bad goal} & \textbf{Good goal} & \textbf{Repair adds} \\
\midrule
``Improve the chapter.'' & Every dimension $\geq 90$ via \texttt{/score-content}.
& Checkable predicate \\
\addlinespace
``Make the model accurate.'' & Held-out AUC $\geq 0.78$ on 2023; no
post-origination features. & Artifact, target, leakage guard \\
\addlinespace
``Get tests passing.'' & \texttt{pytest} 0 \emph{and} coverage $\geq 90\%$;
swallowing exceptions forbidden. & Anti-gaming guard \\
\addlinespace
``Keep improving the portfolio.'' & Max Sharpe s.t.\ turnover $\leq 30\%$; stop
after 8 runs or gain $< 0.05$. & Constraint + stopping rule \\
\bottomrule
\end{tabular}

\medskip
\begin{alertblock}{Each bad goal disables a loop phase}
Unconcrete, uncheckable, or unbounded. \alert{A loop is only ever as good as the
goal you hand it} --- no amount of clever iteration recovers a broken goal.
\end{alertblock}
\end{frame}

% =====================================================================
\section{Agents, Skills, and Hooks}
% =====================================================================

\begin{frame}{Agents as Specialized Iterators --- Five Roles}

Useful agents are not chatbots-with-tools; they are \alert{specialized
iterators}, each doing one transformation on project state and handing off.

\medskip
\begin{itemize}
  \item \textbf{Worker} --- produces/transforms the artifact (the only mandatory role).
  \item \textbf{Critic} --- structured judgment without rewriting; the
        natural-language \emph{gradient}.
  \item \textbf{Validator} --- objective deterministic pass/fail; the role most
        often omitted and most missed.
  \item \textbf{Planner} --- decomposes a goal into ordered subtasks.
  \item \textbf{Refactorer} --- improves structure, preserves behavior.
\end{itemize}

\medskip
\begin{bigpicture}
\textbf{Critic vs.\ validator, in finance terms:} the critic is the senior
reviewer who says the model \emph{looks} aggressive; the validator is the limit
system that simply \emph{rejects} the trade. A loop of only workers + critics can
drift confidently in a wrong direction.
\end{bigpicture}

\medskip
\textbf{Central discipline:} an agent's value is the \emph{durable change to
state} (files, green tests, a commit) --- not the text it prints. Measure it by
its diff.
\end{frame}

\begin{frame}{Skills: Reusable Procedures (a Skill $\approx$ a Function)}

If agents are the \emph{who}, a skill is the \emph{how}: a named, composable
sequence run the same way every time (\texttt{/draft-chapter},
\texttt{/score-content}, \texttt{/iterate-book-quality}).

\medskip
\begin{block}{The skill--function correspondence}
A \emph{name} you invoke it by $\cdot$ \emph{scoped inputs} $\cdot$ a \emph{single
responsibility} $\cdot$ \emph{repeatable} (same inputs $\to$ same procedure)
$\cdot$ \emph{composable} (one skill calls another: \texttt{/draft-chapter}
$\to$ \texttt{/score-content} $\to$ \texttt{/refine-until-threshold}).
\end{block}

\medskip
\textbf{Two further payoffs.}
\begin{itemize}
  \item \textbf{Freeze conventions:} \texttt{/new-topic} creates
        chapter/lecture/notebook $N$ together --- numbering \emph{cannot} drift.
  \item \textbf{Reduce entropy:} every invocation pushes a long-running project
        \emph{back} toward consistency. The skill library is the project's
        control framework --- its standard operating procedures.
\end{itemize}
\end{frame}

\begin{frame}[fragile]{Hooks: Guardrails That Do Not Depend on the Agent}

A \textbf{hook} runs automatically on an event (pre-tool, post-write,
session-end), \alert{outside the agent's control}. Its defining property: it does
not depend on the agent's cooperation.

\smallskip
\textbf{Pre-action} = pre-trade limit check (can block). \quad
\textbf{Post-action} = post-trade surveillance (validate / record / repair / escalate).

\begin{lstlisting}[language=bash]
#!/usr/bin/env bash
# check-citations.sh -- post-write validation hook
set -euo pipefail
BIB="book/bibliography.bib"; MISSING=0
for f in $(git diff --cached --name-only | grep '\.tex$'); do
  for key in $(grep -oE '\\cite\{[^}]+\}' "$f" \
               | sed -E 's/\\cite\{([^}]+)\}/\1/'); do
    grep -q "{$key," "$BIB" || { echo "Unresolved: $key"; MISSING=$((MISSING+1)); }
  done
done
[ "$MISSING" -gt 0 ] && { echo "FAIL"; exit 1; }; echo "OK"; exit 0
\end{lstlisting}

\smallskip
A hook converts ``is this good?'' (opinion) into ``compiles, citations resolve,
scores $\geq$ threshold'' (a \emph{reproducible bit}). Hang the stopping condition
on the bit, not on the model's confidence.
\end{frame}

\begin{frame}{Trustworthy Validators and Failure Hooks}

\begin{underhood}
\textbf{A passing check is necessary, not sufficient.} The loop optimizes for
$V=\text{pass}$, not the true goal $G$. Unless
$\{x: V(x)=\text{pass}\}\subseteq\{x: G\text{ met}\}$, there exist artifacts that
pass yet fail --- and an optimizing loop will find them (Goodhart).
\end{underhood}

\medskip
\textbf{Make $V$ hard to satisfy without meeting $G$:} outcome-tied not proxy
checks; \emph{diverse} checks (game all at once); validator logic out of the
worker's reach; a hidden release-time holdout. (Model-risk management's playbook.)

\medskip
\begin{block}{Failure hooks --- what fires on a red verdict (escalating)}
\textbf{Block} (refuse) $\cdot$ \textbf{Revert} (undo to last-good) $\cdot$
\textbf{Escalate} (hand to a human) $\cdot$ \textbf{Re-enter} (feed the failure
back as next input). The operational-risk playbook in miniature.
\end{block}
\end{frame}

% =====================================================================
\section{Designing an Iteration Loop}
% =====================================================================

\begin{frame}{The Seven-Step Template}

\begin{enumerate}
  \item \textbf{Read state} --- ground truth from disk + version control first.
  \item \textbf{State the next local goal} --- a checkable target for \emph{this} turn.
  \item \textbf{Propose} a change (reason).
  \item \textbf{Apply} it (act --- write to disk; only now has the world changed).
  \item \textbf{Validate} --- a deterministic check; the agent does \emph{not}
        grade its own work.
  \item \textbf{Interpret feedback} --- direction and magnitude of the error.
  \item \textbf{Decide:} stop / continue / \alert{escalate}.
\end{enumerate}

\medskip
\begin{alertblock}{Two non-negotiables}
Validation (5), interpretation (6), and decision (7) are \emph{separate} ---
collapsing them removes the independent check. And the decision always includes
\textbf{escalate}: a loop that can only succeed or grind forever \emph{will}
grind forever.
\end{alertblock}
\end{frame}

\begin{frame}[fragile]{Pseudocode --- and Its Map onto Claude Code Primitives}

\begin{lstlisting}[language={}]
state <- read_state()            # file reads
goal  <- read_goal()             # checked-in goal file
best  <- evaluate(state, goal)
while not goal.satisfied(best) and not goal.budget_exhausted():
    local  <- choose_objective(state, best)
    change <- propose(state, local)     # REASON (ReAct)
    state2 <- apply(change, state)      # ACT: write to disk
    report <- validate(state2)          # a SKILL: deterministic check
    if report.improves_on(best):
        commit(state2, summarize(change, report))   # audit trail
        state, best <- state2, report.score         # accept
    else:
        state <- rollback(state)                     # reject (Self-Refine)
    if stalled(history) or report.requires_human:
        escalate(state, report); break               # hook-gated
emit_report(state, best, history)                     # hook-gated
\end{lstlisting}

\smallskip
The accept/reject discipline makes the loop \alert{robust to a noisy proposer}: a
bad proposal costs one iteration, not the whole run --- because validation and
version control are deterministic.
\end{frame}

\begin{frame}{State Lives on Disk; Stop on a Budget}

\begin{bigpicture}
\textbf{Design for amnesia.} The agent has no durable memory between turns ---
it is a new employee every morning who reads the files, does a day's work, commits
it, and is replaced overnight. \alert{Write before you forget.}
\end{bigpicture}

\medskip
\textbf{Persists:} files, git history (long-term memory), validation reports, a
human-readable status scratchpad. \textbf{Lost:} chain-of-thought, unsaved
values. Bonus: all-on-disk state makes the loop \emph{resumable}.

\medskip
\begin{block}{Stopping rules --- exactly one must be guaranteed-terminating}
Success $\cdot$ convergence (no gain for $k$ steps) $\cdot$ \alert{budget
exhaustion (mandatory)} $\cdot$ human-review-required. Convergence/escalation make
the loop \emph{smart}; the budget makes it \emph{safe}.
\end{block}

\medskip
\textbf{Economics:} marginal value shrinks near the goal while cost stays flat ---
stop where marginal value crosses marginal cost (latency binds interactively;
dollars bind overnight).
\end{frame}

% =====================================================================
\section{Example 1: Derivative-Free Root Finding}
% =====================================================================

\begin{frame}{Example 1: A Loop That Improves Using Feedback Alone}

\textbf{Find} $x$ with $f(x)=x^3-x-2\approx 0$ to tolerance $|f(x)|<10^{-3}$
(single root near $x\approx 1.52$).

\medskip
\alert{The point is not numerical analysis.} The agent may \emph{evaluate} $f$
anywhere but is \textbf{forbidden derivatives} --- no $f'(x)=3x^2-1$, no Newton
update. It reads only the sign and size of the residual. This is the whole
machinery, stripped bare.

\medskip
\begin{bigpicture}
\textbf{Finance analogue:} this is mechanically an \emph{IRR} or
\emph{implied-volatility} solver. For an IRR, seek $r$ with
$\mathrm{NPV}(r)=\sum_t c_t/(1+r)^t=0$ --- no closed form, so evaluate at trial
rates and move by the sign (too-high NPV $\Rightarrow$ rate too low). Swap $f$ and
the loop is unchanged.
\end{bigpicture}

\medskip
\textbf{Three artifacts, none containing an algorithm:} a \emph{goal file}
($|f|<0.001$, $\leq 12$ iters), a \emph{skill} (\texttt{python evaluator.py \$x |
tee -a results.log}), and a \emph{stop hook} that vetoes ``success'' unless the
last $|f(x)|<10^{-3}$.
\end{frame}

\begin{frame}{Iteration Trace from a Deliberately Poor Start ($x=0$)}

\begin{table}
\centering
\begin{tabular}{@{}rrrl@{}}
\toprule
Iter & $x$ & $f(x)$ & Agent's reasoning \\
\midrule
1 & $0$      & $-2$        & Negative; root to the right. Jump up. \\
2 & $2$      & $4$         & Positive; root bracketed in $(0,2)$. \\
3 & $1.5$    & $-0.125$    & Negative but small; root just above $1.5$. \\
4 & $1.6$    & $0.496$     & Positive; bracket $(1.5,1.6)$. \\
5 & $1.52$   & $-0.008192$ & Tiny negative; nudge up. \\
6 & $1.522$  & $\approx 0.004$ & Overshot; back off a touch. \\
7 & $1.5214$ & $\approx 0$ & $|f(x)|<0.001$: hook passes. \\
\bottomrule
\end{tabular}
\end{table}

\smallskip
The \alert{sign change} between $0$ and $2$ is the engine: a continuous function
going $-\to+$ crosses zero between. The magnitudes ($|f(1.5)|\ll|f(1.6)|$) place
the root near $1.5$ (false-position reading). \textbf{Declaring success at
iteration~5 would be vetoed by the hook.}

\medskip
\begin{block}{Why the inefficiency is the point}
The zig-zag past the root means the loop reached tolerance \emph{without an
algorithm} --- so the same loop works when the residual is a failing test, a
style violation, or a reviewer's objection.
\end{block}
\end{frame}

% =====================================================================
\section{Example 2: Credit-Risk Stress-Testing Loop}
% =====================================================================

\begin{frame}{Example 2: A Defensible Credit-Committee Memo}

\textbf{Meridian Components} (fictional; all figures illustrative): a mid-market
manufacturer seeking a \$40\,m revolver. Looks comfortable (growth, rising EBITDA,
leverage in appetite) --- but cyclical end markets, illiquid inventory,
concentrated receivables.

\medskip
\begin{alertblock}{The goal is \emph{not} ``estimate the PD''}
A point estimate cannot be \emph{wrong} --- nothing fails against it. The goal is
a \emph{memo a committee would accept}: that acceptance is checkable, so the loop
runs until the checks pass.
\end{alertblock}

\medskip
\textbf{The expected-loss decomposition} (reconciled in every column):
\[
  \mathrm{EL} \;=\; \mathrm{PD}\times\mathrm{LGD}\times\mathrm{EAD}.
\]
PD = 1-yr default prob.; $\mathrm{LGD}=1-\text{recovery}$; EAD = drawn balance
plus drawdown. The memo must give PD, LGD, EAD, EL (\$ and bps), drivers, and a
recommendation --- \emph{lend / modify / decline}. A ``low risk'' phrase is gated
on both base \emph{and} stressed EL clearing thresholds.
\end{frame}

\begin{frame}{The Combined Shock and the Nine-Agent Ensemble}

\begin{columns}[T]
\begin{column}{0.46\textwidth}
\textbf{Combined macro shock} (parameter table, not prose):
\smallskip
\begin{tabular}{@{}lrr@{}}
\toprule
Parameter & Adv. & Sev. \\
\midrule
Revenue & $-20\%$ & $-30\%$ \\
Gross margin & $-300$ & $-450$ bps \\
Rate & $+250$ & $+350$ bps \\
DSO & $+25$ & $+40$ d \\
Inventory days & $+20$ & $+35$ d \\
ICR covenant & \multicolumn{2}{c}{$\to 2.5\times$} \\
\bottomrule
\end{tabular}
\smallskip\\
\alert{Combined}, because the drivers are correlated; breaches trigger only when
several coincide.
\end{column}
\begin{column}{0.52\textwidth}
\textbf{Nine role-scoped agents} (each output $\to$ next input):
\smallskip
\begin{enumerate}\footnotesize
  \item Data-ingestion
  \item Accounting-quality (normalize)
  \item Feature-engineering
  \item PD-model \hfill\textit{(leakage hook)}
  \item Stress-scenario
  \item Cash-flow (covenants, liquidity)
  \item LGD (collateral, recovery)
  \item \alert{Skeptic} --- returns to ANY agent
  \item Memo-writing \hfill\textit{(skeptic-gated)}
\end{enumerate}
\smallskip
Two \emph{feedback edges} (skeptic + failing hooks) turn a pipeline into a loop.
\end{column}
\end{columns}

\medskip
\begin{block}{The two load-bearing hooks}
\textbf{Leakage check} (the easy path to a beautiful backtest is to peek at the
answer) and the \textbf{``low risk'' gate} (language is a control surface ---
earned by arithmetic, not chosen by the writer).
\end{block}
\end{frame}

\begin{frame}{Eleven Iterations, One Exhibit (EAD $=$ \$40\,m)}

\begin{columns}[T]
\begin{column}{0.40\textwidth}
\footnotesize
\textbf{The arc:}
\begin{enumerate}\footnotesize
  \item Naive PD $\approx 1.2\%$
  \item Add-backs $\Rightarrow$ lev.\ $3.0\to3.8\times$
  \item Richer features
  \item \alert{Leakage hook rejects} PD $0.9\%$
  \item Clean rebuild $\to$ PD $2.5\%$
  \item Adverse stress $\to$ PD $9\%$
  \item Covenant breach; $\sim$\$6\,m hole
  \item LGD $40\%\to55\%$
  \item Skeptic demands severe + peers
  \item Severe: PD $13\%$, LGD $60\%$
  \item Conditional approval
\end{enumerate}
\end{column}
\begin{column}{0.60\textwidth}
\begin{table}
\centering\small
\begin{tabular}{@{}lrr@{}}
\toprule
Metric & Base & Stressed \\
\midrule
PD & $2.5\%$ & $9.0\%$ \\
LGD & $40\%$ & $55\%$ \\
EAD & \$40.0\,m & \$40.0\,m \\
EL (\$) & \$0.40\,m & \$1.98\,m \\
EL (bps) & 100 & 495 \\
Covenant breach? & No & Yes \\
Liquidity short. & --- & $\approx$\$6.0\,m \\
\bottomrule
\end{tabular}
\end{table}
\smallskip
\footnotesize
$\mathrm{EL}_{\text{base}}=0.025\cdot0.40\cdot40=\$0.40$\,m;\;
$\mathrm{EL}_{\text{str}}=0.09\cdot0.55\cdot40=\$1.98$\,m.
\smallskip\\
\textbf{Action:} conditional --- \$30\,m, tighter covenants, more collateral,
repriced.
\end{column}
\end{columns}

\medskip
\begin{alertblock}{The story the loop found, the naive memo missed}
Sound \emph{in expectation}, fragile \emph{in the tail}: EL nearly quintuples, and
the damage is \emph{structural} (breach + liquidity hole). A 1.2\% point estimate
hid all of it. A \emph{recommendation}, not a decision --- a human signs.
\end{alertblock}
\end{frame}

% =====================================================================
\section{Oversight, Failure Modes, and Patterns}
% =====================================================================

\begin{frame}{Human Oversight, Sign-Off, and Audit Trails}

\textbf{What does the signature add?} Hooks answer ``did it pass the checks?'',
never ``who is \emph{answerable}?'' The loop manufactures analysis, not
accountability or legal standing.

\medskip
\begin{itemize}
  \item \textbf{Escalation triggers} (stated in the goal file, not improvised):
        materiality threshold crossed, low confidence / out-of-range inputs,
        agents that cannot reconcile, a hook firing repeatedly, irreversible
        action. \emph{Human-review-required} is a terminal state.
  \item \textbf{Sign-off} (SR 11-7) = an event with time, identity, scope --- not
        a property of the artifact. The skeptic is internal challenge, \alert{not}
        independent validation. Carry an \emph{unsigned sign-off block} the loop
        cannot complete.
  \item \textbf{Audit trail:} version control $+$ structured run log $+$
        content-addressed inputs $\Rightarrow$ an auditor can \emph{regenerate}
        every number. ``Reproducible'' means the derivation is recoverable and
        therefore challengeable.
\end{itemize}
\end{frame}

\begin{frame}{Failure Modes --- and the Control That Catches Each}

\begin{table}
\centering\small
\begin{tabular}{@{}p{0.24\linewidth}p{0.34\linewidth}p{0.34\linewidth}@{}}
\toprule
\textbf{Failure mode} & \textbf{Warning sign} & \textbf{Control} \\
\midrule
Infinite loop & Iters climb, no change; oscillation & Hard budget + cost ceiling \\
\addlinespace
Fake convergence & ``Done'' but metric plateaued below goal & Terminate on
\emph{external} validator, not self-report \\
\addlinespace
Goal drift & Later iters optimize a different objective & File-based goal re-read
each iteration \\
\addlinespace
Premature success & ``Low risk'' before the number is met & Tolerance hook \\
\addlinespace
Goodhart overfit & One check passes; OOS lags & Diverse + held-out + adversarial \\
\addlinespace
Excessive delegation & Agents for deterministic subtasks & Functions for
closed-form; when-not-to-loop test \\
\bottomrule
\end{tabular}
\end{table}

\medskip
\begin{alertblock}{The meta-point}
Every control is a mechanism introduced earlier. \alert{The failure modes are the
predictable consequences of omitting machinery the lecture already argued for} ---
build the loop right and you have defended against most of them by construction.
\end{alertblock}
\end{frame}

\begin{frame}{Five Design Patterns}

\begin{itemize}
  \item \textbf{Planner--executor--critic} --- planning + execution + independent
        evaluation; for tasks where the next step depends on earlier results.
  \item \textbf{Generate--test--revise} --- the tightest loop; default when the
        goal is ``make this check pass.'' Only as good as its test.
  \item \textbf{Scaffold--draft--audit--repair} --- for subjective, rubric-scored
        goals. \alert{This is the pattern behind this very book.}
  \item \textbf{Data-pipeline: build--validate--document} --- the artifact is a
        process; \emph{lineage} must be trustworthy.
  \item \textbf{Credit-model: build--stress-test--red-team--finalize} ---
        Example~2 generalized; stress-test and red-team are \emph{mandatory}
        stages (SR 11-7 governance encoded as a pattern).
\end{itemize}

\medskip
\begin{block}{Choosing --- and composing}
Pick by the goal's properties (machine-checkable vs.\ judgment; number vs.\
artifact; stakes; reproducibility). Real workflows \emph{nest} patterns. What
never changes: the feedback must be \alert{informative}.
\end{block}
\end{frame}

% =====================================================================
\section{Summary and Takeaways}
% =====================================================================

\begin{frame}{Lecture 17 Summary}

\textbf{The shift:} from issuing instructions and accepting the reply, to
\alert{designing a goal-directed loop} whose output is observed, evaluated, and
revised against an explicit objective.

\medskip
\textbf{Three building blocks:} \emph{agents} specialize the loop (\emph{who}),
\emph{skills} reuse it (\emph{how}), \emph{hooks} enforce it (\emph{what it may
never violate}).

\medskip
\textbf{Three controls:} goals (the setpoint), validators (honest feedback),
stopping rules (terminate safely).

\medskip
\begin{alertblock}{The one idea to internalize}
A loop adds value \emph{exactly to the degree each iteration's feedback informs
the next action.} In finance the goal is never to produce an answer, but to
produce an answer someone can stand behind --- achieved not by a better prompt
but by a \alert{better loop}.
\end{alertblock}

\medskip
\textbf{Next (practical session):} hand-run the root-finding loop, write a
tolerance success-gate in bash, and decide lend/no-lend from base vs.\ stressed EL.
\end{frame}

% =====================================================================
\appendixstart{Definitions, Literature, and Reference Material}
% =====================================================================

\begin{frame}[noframenumbering]{Reference: The Agentic-AI Lineage}

\begin{block}{LLM side --- the loop generalizes a family of methods}
\textbf{ReAct} (Yao 2022): interleave reason + act = Observe--Propose--Act.
\textbf{Reflexion} (Shinn 2023): verbalized Evaluate--Revise.
\textbf{Self-Refine} (Madaan 2023): single-model generate--critique--revise.
\textbf{Chain-of-thought} (Wei 2022); \textbf{Toolformer} (Schick 2023);
\textbf{AutoGen} (Wu 2023): split phases across role-typed agents.
\end{block}

\begin{block}{Engineering side --- the loop predates LLMs by decades}
Boehm's \textbf{spiral model} (1988): risk-driven cycles replace waterfall.
Deming's \textbf{PDCA}: Plan--Do--Check--Act. The convergence is no coincidence:
it is what any system producing correct output under uncertainty looks like once
it closes its feedback loop.
\end{block}

\medskip
\footnotesize\textbf{Control-theoretic reading:} one-shot $=$ open-loop;
goal-directed $=$ feedback control. Error $e_{t+1}=G-s_{t+1}$; instability when
gain is too high (oscillation); steady-state error when the success predicate is
biased.
\end{frame}

\begin{frame}[noframenumbering]{Reference: The Meridian Acceptance Criteria}

\begin{table}
\centering\small
\begin{tabular}{@{}lll@{}}
\toprule
Deliverable & Acceptance criterion & Checked by \\
\midrule
PD (base \& stressed) & OOS-validated; horizon stated & Hook + skeptic \\
LGD & Tied to collateral; recovery logic & Skeptic \\
EAD & Drawdown assumption stated & Hook \\
Expected loss & $\mathrm{EL}=\mathrm{PD}\cdot\mathrm{LGD}\cdot\mathrm{EAD}$ reconciles & Hook \\
Risk drivers & Each traced to a source statement & Hook \\
Scenarios & Base / adverse / severe, parameterized & Hook \\
Recommendation & lend / modify / decline, justified & Skeptic \\
``Low risk'' & Only if base \emph{and} stressed EL below thresholds & Hook \\
\bottomrule
\end{tabular}
\end{table}

\medskip
\begin{block}{The test for ``goal'' vs.\ ``task''}
Ask what it would mean for the work to be \emph{wrong}. ``Estimate the PD'' cannot
be wrong. ``A memo whose stressed EL reconciles, whose every ratio ties to source,
and which the skeptic cannot break'' can be wrong in a dozen concrete ways --- and
each is a hook waiting to fire.
\end{block}
\end{frame}

\begin{frame}[fragile,noframenumbering]{Reference: The Root-Finding Stop Hook}

\begin{lstlisting}[language=bash]
#!/usr/bin/env bash
# stop-hook: veto "success" unless the latest |f(x)| < 0.001
last=$(grep -o 'f(x)=[-0-9.eE]*' results.log | tail -n1 | cut -d= -f2)
abs=$(python -c "print(abs(float('$last')))")
ok=$(python -c "print(1 if abs(float('$last')) < 0.001 else 0)")
if [ "$ok" -ne 1 ]; then
  echo "BLOCKED: |f(x)|=$abs >= 0.001 -- not converged" >&2
  exit 1
fi
echo "OK: |f(x)|=$abs < 0.001"
exit 0
\end{lstlisting}

\smallskip
The harness treats a nonzero exit from a stop hook as a \emph{veto}: the agent
literally cannot end the loop claiming success while the residual is too large.
The goal says \emph{what} convergence means; the skill says \emph{how} to measure;
this hook \emph{enforces} it.
\end{frame}

\end{document}
